
\documentclass{article} %210 mm × 297 mm
\usepackage[utf8]{inputenc}
\usepackage[a4paper,left=1.5cm, right=1.5cm, top=2cm, bottom=2cm]{geometry}
\usepackage{graphicx}
\usepackage{systeme}          % for Solution 3
\usepackage{array}  
%\usepackage{blindtext}
\usepackage{multirow}
\usepackage[normalem]{ulem}
\usepackage{mhchem}
\usepackage{url}
\usepackage{subfigure} 
\usepackage{amsfonts,latexsym} % para tener disponibilidad de diversos simbolos
\usepackage{enumerate}
\usepackage{longdivision}
\usepackage{polynom}
%\usepackage{booktabs}
\usepackage{float}
%\usepackage{threeparttable}
\usepackage{array,colortbl}
%\usepackage{ifpdf}
%\usepackage{rotating}
\usepackage{stfloats}
\usepackage{url}
\usepackage{listings}
\usepackage{fancyhdr}
\usepackage{biblatex}
\usepackage{color}
\addbibresource{sources.bib}
%\usepackage{hyperref}
%\usepackage{wrapfig}
\usepackage{array}
%\usepackage{multirow}
%\usepackage{adjustbox}
%\usepackage{nccmath}
\usepackage[dvipsnames]{xcolor}
\usepackage{tcolorbox}
\newtcolorbox{mybox}{colback=white,
colframe=black}
\newcommand{\titulo}{Gedächtnisprotokoll Zweitklausur}
\newcommand{\nombre}{Diskrete Mathematik}
\newcommand{\FCE}{\textbf{WiSe 2024/2025}}

 

\usepackage{fancyhdr}
\pagestyle{fancy}
\fancyhead{}
\fancyfoot{}
\fancyhead[LO]{\small\nombre}
\fancyhead[RE]{\small\titulo}
\fancyhead[LO]{\small\FCE}
\fancyhead[RO]{\small }
\fancyfoot[RO,LE] {\thepage}
\usepackage[onehalfspacing]{setspace}

\setcounter{page}{1} %Número de la página, la editorial lo cambiará
\begin{document}
\begin{tabular}{p{12cm}}
{{\Large\bf\titulo}}\\
\vspace{0.2cm}\\
 {\large\nombre}\\
 \vspace{0.1cm}
\end{tabular} 



\section{Permutationen}
(Permutation nicht genau so...aber mit den wiederholten Zahlen eben)
\begin{equation*}
    S_8, \quad \pi = (6,1,4)(3,1)(2,4,7)
\end{equation*}

\begin{enumerate}
    \item Bestimmt alle $\pi(x)$
    \item Schreibe $\pi$ als Menge disjunkter Zykel
    \item finde ein $\sigma \neq \pi$ das mit $\pi$ konjugiert ist 
    \item Bestimme den Typ von $\sigma$ und $\pi$ und erkläre wie die Typen zusammen hängen 
\end{enumerate}

\section{Äquivalenzklassen oder so?}
\begin{enumerate}
    \item Inverses von $[463]$ bestimmen in $Z / 151 Z$
    \item Polynom mit Grad $\leq 6$ finden dass gleiche Restklasse hat wie ein vorgegebenes Polynom 
    \begin{itemize}
        \item Vorgegebenes Polynom war $x^{11} + x^8 + x^6 + x^5 + 1$ oder so (auf jeden Fall Grad 11) 
        \item $k[x]$
        \item Da war noch was mit $Z / 151 Z \backslash (x^7 + iwas)$...
    \end{itemize}
\end{enumerate}

\section{Designs}
\begin{enumerate}
    \item Finde/Erstelle in Design mit (4,2,2) (da stand NICHT t-Design oder so) und begründe, warum man mit (4,2,2) ein gültiges Design hat 
    \item Bestimme eine Differenzenfamilie in/über(?) $Z/11Z$ mit 5 Elementen und erstelle daraus ein Design
\end{enumerate}

\section{Codes}
Kontrollmatrix gegeben
\begin{enumerate}
    \item Finde alle Codewörter des Codes (aus der Kontrollmatrix)
    \item Wie viele Fehler kann der Code korrigieren? Begründe
    \item Du empfängst ein Codewort 01...., korrigiere es unter der Annahme, dass es nur einen Fehler enthält
\end{enumerate}

\section{Rekursionsgleichungen}

\begin{equation*}
    u_n = u_{n-1} + 2 \cdot 5^n
\end{equation*}
Anfangswert $u_0 = -3$

\begin{enumerate}
    \item Berechne $u_1$ und $u_2$
    \item für die Folge $u_1, u_2, u_3$ 
    \item Lösung der Rekursion finden 
\end{enumerate}

Folgende Formeln (so oder so ähnlich) waren als "vielleicht hilfreich" gegeben: 
\begin{equation*}
    \sum_{n \geq 0} x^n = \frac{1}{1-x}, \quad \sum_{n \geq 0} \underbrace{(n-1)}_{oder (x-1?)}x^n = \frac{1}{(1-x)^2}
\end{equation*}

\end{document}